\chapter*{}
\noindent{\textbf{RESUMO}}

Este trabalho apresenta o desenvolvimento de um jogo 2D do gênero \textit{cozy}. A motivação surgiu da necessidade de explorar um tipo de jogo ainda pouco abordado na literatura, que valoriza mecânicas simples e ambientação convidativa. O objetivo é criar um jogo que una simplicidade, estética aconchegante e funcionalidades que promovam conforto e imersão. Para isso, foi realizado um estudo sobre a história e a classificação dos jogos, além da análise de motores gráficos, resultando na escolha da \textit{engine Godot}, por sua flexibilidade, facilidade de uso e ampla documentação. A metodologia adotada envolveu etapas como definição do estilo visual e narrativo, seleção e adaptação de \textit{assets} e desenvolvimento das funcionalidades centrais. Como resultado, foi desenvolvido um jogo que cumpre os objetivos propostos, apresentando uma estética consistente, jogabilidade fluida e demonstrando o potencial da \textit{Godot} para projetos desse tipo. Espera-se que este trabalho contribua para o fortalecimento da produção independente de jogos \textit{cozy} e sirva de base para estudos futuros na área.

\vspace{\onelineskip}

% todas em letras minúsculas, separadas por ponto e vírgula (;)
\noindent{\textbf{Palavras-chave}: \textit{Cozy Game}. Godot. 2D. Jogos.}